\chapter{Congestion control}
The congestion control is used limit the transmission rate when the network is
diagnosted to suffer from congestion. Several factors can be used to detected
the congestion. Traditional TCP variants most often use packet losses as an
indicator of congestion, some others measure the latencies and deduce the
increased queueing delay as first signs of congestion. We can also intelligent
middleboxes to implement ECN functionality, which adds additional signalling to
inform end hosts about the congestion. We chose to use packet losses as it is
the one which works reasonable well without large scale testing and analysis
(which is most likely required to fine tune the delay based congestion control mechanishms) or
without any additional mechanisms from the network.

The setting we have in this assignment assumes that we cannot control the sensor nodes
so nothing can done to improve the behaviour under congestion on that link. So
we study only the congestion on the link between the server and the clients.
Basically the server is the network intensive part of the protocol
implementation as it creates
the most of the protocol messages under normal conditions.

The protocol has three communication patterns the server can use to measure and
deduce that a loss has occured. The first one is a missing update
acknowledgement when reliable subscription is used. The another one is to rely
on negative acknowledgements sent by the client who can deduce a loss from
missing sequence numbers and send the NACK message even if the communication
flow is not set up to be reliable.

When a loss is detected the server must either \textit{skip messages} or
\textit{throttle down} its sending rate. The first one is a more simple solution
which causes the server simply to skip some updates from being delivered. The
latter option will cause the server to send the updates at slower pace. Because
we use weak reliability guarantees this also means that if the sensors publish
updates at a pace greater than the server can send the updates forward, some of
these updates will be discarded.


\subsection{Skipping messages}
Because flows without reliability guarantees are not communicating back to
server server has no means of reacting quickly to congestion situations by
default. One simple idea is to add negative acknowledgements into protocol
design. This is a mechanisms which is generally considered quite inferior as
more messages are added into already congested network (e.g. see why ICMP source
quench is deprecated).

When a client detects a cap in the sequence number flow the client will generate
a special NACK packet and sends it to server. Server receiving the packet will
deduce that a link is congested and will simply skip the next message send to
client. When a message is skipped the server will reset the lost message state
and the next message after the skipped message will be delivered normally.

Message skipping is happening on the link between the server and the client
shared by multiple subscriptions.

\subsection{Throttling down the sending rate}
When a server detects congestion it can also implement more sophisticated
methods to handle the situation like to throttle down its sending rate. This
means that server will schedule the transmission of update later under the
congestion. With our reliability guarantees this also means that a scheduled
message might be replaced by a newer message which means that it's actually
never transmitted.


\subsection{After thoughts about about our congestion control design}
Adding congestion control on top if this kind of application turned out to be
quite taunting task. First of all we have a possibility of two kind of clients
first of which are basically only receiving updates and never acknowledging
these back to server which however require a method of informing server about
the congestion. Another problem is the weak-reliability guarantee which adds its
own complexity due to fact that not all messages are always even supposed to be
received (e.g. in case of update bursts when a new message quickly replaces the
old one).

